\documentclass[11pt]{article}


\setlength{\parindent}{0pt}
\usepackage{xltxtra}
\usepackage{hyperref}
\hypersetup{hidelinks}
\usepackage{url}
\urlstyle{tt}
\usepackage{xcolor}
\definecolor{CVBlue}{RGB}{23,110,191}
\usepackage{calc}
\usepackage{graphicx}
\usepackage{tikz}
\usetikzlibrary{calc}
\usepackage{fontspec}
\usepackage{xeCJK}
\usepackage{enumitem}
\CJKsetecglue{} %% 取消中文与数字之间的间隙


%% 主文档字体设置
\setmainfont[
    Path = fonts/Main/,
    Extension = .otf,
    BoldFont = texgyretermes-bold.otf, % 加粗字体
]{texgyretermes-regular.otf} % 正文字体

% 中文字体设置
\setCJKmainfont[
    Path = fonts/hansans/,
    Extension = .ttf,
    BoldFont = NotoSansSC-Bold.ttf, % 加粗字体
]{NotoSansSC-Regular.otf} % 正文字体


\usepackage{relsize}
\usepackage{xspace}

% 使用 fontawesome（部分图标）
\usepackage{fontawesome} 

% A4纸，上下左右边距
\usepackage[
    a4paper,
    left=1.2cm,
    right=1.2cm,
    top=1.5cm,
    bottom=1cm,
    nohead
]{geometry}

\renewcommand{\baselinestretch}{1.5} % 行间距设为1.5

\usepackage{titlesec}
\usepackage{enumitem}
\setlist{noitemsep} % 取消列表项间的额外间距
%\setlist{nosep} % 取消所有垂直间距
\setlist[itemize]{topsep=0.25em, leftmargin=*}
\setlist[enumerate]{topsep=0.25em, leftmargin=*}

% --- 用于控制【不同项目之间】的垂直距离 ---
\newlength{\interProjectSpacing}
\setlength{\interProjectSpacing}{0.5em} % <--- 在此调整项目之间的距离
\newcommand{\projectsep}{\vspace{\interProjectSpacing}}

% --- 用于控制【项目标题】与下方【项目描述】的距离 ---
\newlength{\intraProjectTitleSep}
\setlength{\intraProjectTitleSep}{0.2em} % <--- 在此调整标题和描述的距离
\newcommand{\titlebreak}{\\[\intraProjectTitleSep]}

% --- 用于控制【项目描述】与下方【要点列表】的距离 ---
\newlength{\intraProjectListTopSep}
\setlength{\intraProjectListTopSep}{0.2em} % <--- 在此调整描述和列表的距离

% =======================================================================


\titleformat{\section}         % 定制 \section 命令 
{\large\bfseries\raggedright} % 将 section 标题设置为大号、粗体且左对齐
{}{0em}                      % 可用于为所有 section 添加前缀（如“章节...”）
{}                           % 可用于在标题前插入代码
[{\color{CVBlue}\titlerule}]  % 在标题后插入一条横线
\titlespacing*{\section}{0cm}{*1.6}{*1.2}



\begin{document}
\pagenumbering{gobble}
  %%%% 利用tikz来定位学校Logo，这里只在第一页显示
  \begin{tikzpicture}[remember picture, overlay] 
    \node[anchor = north west] at ($(current page.north west)+(0.5cm,+1.0cm)$) {\includegraphics[height=6cm]{zju.png}};
  \end{tikzpicture}%
\centerline{\LARGE\bfseries{杨晓菡}} 

\centerline{\normalsize{\faPhone\ 19818528261 \quad \faEnvelopeO\ \href{mailto:3230100000@zju.edu.cn}{3240106073@zju.edu.cn}}} 
    
\section{\makebox[\widthof{\faGraduationCap}][c]{\color{CVBlue}\faGraduationCap}\ 教育背景}    
\textbf{浙江大学} \hfill 2024.9 -- 至今\\[0.5em] % 标题和正文间加一点距离
信息工程\quad 大二 
\begin{itemize}[nosep]
    \item 相关课程：《电子电路基础》《电子电路设计实验》《信号与系统》《电磁场与电磁波》《数字系统设计》《人工智能》《概率论与数理统计》
    \item 主修均绩：4.06/5.0
\end{itemize}

\section{\makebox[\widthof{\faUsers}][c]{\color{CVBlue}\faUsers}\ 个人能力}

% --- 第一个项目 ---
% 将标题行末尾的 \\ 替换为 \titlebreak 命令
\textbf{C语言，Verilog HDL，Python}  \titlebreak

\begin{itemize}[nosep, topsep=\intraProjectListTopSep]
    \item \textbf{大一上学期}：系统学习C语言，建立良好的编程思维。
    \item \textbf{大二}：学习Verilog HDL和Python，会用LaTeX排版学术文档，能使用MATLAB进行数据分析与仿真。
   
\end{itemize}

% 使用 \projectsep 命令来分隔两个项目
\projectsep

% --- 第二个项目 ---
\textbf{硬件基础与实验能力}  \titlebreak

\begin{itemize}[nosep, topsep=\intraProjectListTopSep]
    \item \textbf{焊接与制作}：具备直插件及贴片器件手工焊接经验，能完成DIP封装芯片、排针等元器件的焊接与拆换，有学习精细焊接的意愿。
    \item \textbf{仪器操作}：具备数字示波器基础操作能力，能完成波形捕获、电压/周期测量、通道设置等基本功能；了解信号发生器、直流电源的常规使用方法；实验过程中注重规范操作与安全习惯。
    
\end{itemize}
\projectsep
\textbf{工程习惯}  \titlebreak

\begin{itemize}[nosep, topsep=\intraProjectListTopSep]
   \item 实验前制定简单测试计划，实验后整理仪器与工作台，能按电路图进行基础故障排查。

\end{itemize}
\projectsep
\textbf{学习态度与科研素养}  \titlebreak
\begin{itemize}[nosep, topsep=\intraProjectListTopSep]
    \item \textbf{自主学习}：面对新工具或知识，能通过官方文档、技术博客自学，能够熟练使用AI工具进行技能学习。
    \item \textbf{责任心}：对待任务坚持按时交付，认真对待任务。
    \item \textbf{时间投入}:每周可稳定投入10-12小时参与SRTP，愿意从基础任务做起。

    
\end{itemize}

    
\section{\makebox[\widthof{\faInfo}][c]{\color{CVBlue}\faInfo}\ 其他}
\begin{itemize}[parsep=0.5ex]
    \item \textbf{英语水平：} CET-4（666分）, CET-6（571分）,能较为流畅地阅读外语文献，理解论文核心内容与实验方法
\end{itemize}
\end{document}
